\section{Related Work}
\label{sec:related_work}

\subsection{Plagiarism Detection and Novelty Assessment}

Assessing novelty is often confused with plagiarism checking, but the two tasks differ substantially in goal and method. Commercial tools such as Turnitin Similarity \citep{turnitin} and iThenticate \citep{ithenticate} measure surface text overlap and flag missing or inconsistent citations and attributions. They do not assess novelty at the level of ideas, methodology, or contribution; their outputs function as similarity indicators rather than conclusions about conceptual novelty.

When evaluating a research idea, conceptual overlap and textual overlap must be distinguished. Two papers can describe the same underlying idea with little shared wording, while a genuinely new study may still need to reuse common terms such as standard metrics and widely accepted definitions. Similarity reports therefore support academic integrity checks but do not directly answer whether an underlying idea is original.

\subsection{Literature Search and Discovery Tools}

A wide range of tools has emerged to support finding relevant literature and mapping citations, which is a separate goal from quantitative novelty scoring. Semantic Scholar \citep{semanticScholar} applies machine-learning methods to support literature search and the exploration of publication metadata. Connected Papers \citep{connectedPapers} and Research Rabbit \citep{researchRabbit} use graph-based methods to map relationships among related studies. Elicit \citep{elicit} reduces the manual effort in systematic literature reviews by automating parts of screening and data extraction. Scite \citep{scite} uses Smart Citations to label cited works as supporting, contrasting, or mentioning, helping readers judge the strength and direction of evidence. While these systems improve screening and triage, they typically do not provide a clear, rule-based originality score with sentence-level tracking that can be checked against a specific user hypothesis.

\subsection{Literature-Grounded Idea Novelty Systems}

Relying solely on an LLM to judge novelty is risky, because the model can present conclusions confidently even when those claims cannot be checked against verifiable evidence. Retrieval-augmented generation (RAG) is a widely used remedy: it combines text generation with retrieval from an external corpus to improve factual accuracy and to record the sources used \citep{lewis2020rag,gao2024ragsurvey}. The broader notion of novelty itself has earlier roots in information retrieval, where the TREC Novelty Track \citep{harman2002trec} formalised novelty as the sentence-level identification of new information in a document stream. Although the target object differs from scientific idea novelty, the conceptual move from document-level relevance to sentence-level novelty remains influential and motivates the sentence-level overlap reports produced by Hypothetica.

Several recent systems target literature-grounded idea novelty more directly. \citet{shahid2025} propose the Idea Novelty Checker, a retrieve-then-rerank pipeline that evaluates the novelty of a scientific idea using retrieved literature, expert-annotated examples, and a facet-based operational definition of novelty. The output is binary, novel or not novel, and the judgement depends on whether the idea differs from prior work in core facets such as purpose, mechanism, or evaluation. The work demonstrates that novelty assessment quality depends strongly on retrieval quality and on having a clearer operational notion of novelty than generic semantic similarity. In contrast, Hypothetica uses a multi-agent design with separate stages for retrieval, paper-by-paper analysis, and synthesis, draws candidates from multiple corpora through a unified adapter layer rather than from a single literature index, and replaces a binary verdict with a four-dimensional Likert-based assessment whose criterion weights are calibrated against an arXiv benchmark.

\citet{lin2025schnovel} take a complementary, benchmark-first angle. They construct SchNovel, a 15,000-pair benchmark of academic papers across six fields, and propose RAG-Novelty, a retrieval-augmented pipeline that assigns a 0-to-10 novelty score to a query paper after comparing it against top-K retrieved prior work. Their ground truth is operationalised by treating the newer paper in each pair as the more novel one, which scales well but is acknowledged as a debatable proxy across domains. Hypothetica differs in three ways. It scores user-supplied research ideas rather than already-published papers, it replaces a single 0-to-10 axis with four independently-scored criteria with empirically calibrated weights, and it grounds each score in sentence-level evidence rather than abstract-only retrieval.

A second direction broadens the task beyond binary novelty judgement. ScholarEval \citep{scholareval} evaluates research ideas on two criteria, soundness and contribution, using a multi-stage pipeline that retrieves method-level and dimension-level evidence from literature before synthesising strengths, weaknesses, contradictions, and suggestions. ScholarEval moves toward richer reviewer-style feedback rather than a single novelty verdict. Compared to ScholarEval, Hypothetica is narrower in scope, focused specifically on originality assessment, but it similarly emphasises structured, evidence-grounded analysis instead of a single opaque judgement. Hypothetica also broadens the notion of relevant literature to include patent records and public code repositories alongside academic works.

RINoBench \citep{rinobench} formalises research idea novelty assessment as the joint prediction of a rubric-based 1-to-5 novelty score and a grounded textual justification. It is benchmark-oriented rather than system-oriented: it argues that novelty-assessment systems should be evaluated not only on the score they assign, but also on whether their explanations align with grounded human reasoning. This is particularly relevant for systems such as Hypothetica, because plausible reasoning and accurate scoring do not always coincide.

OpenNovelty \citep{opennovelty} focuses on verifiable scholarly novelty assessment. It extracts a paper's core task and claimed contributions, retrieves prior work through expanded semantic queries, organises related work into a hierarchical taxonomy, and performs claim-level comparisons against retrieved papers. Refutation-style judgements are retained only when supported by verified evidence quotes from the source texts. This evidence-bound design is consistent with the broader literature on scientific claim verification, where systems trained on benchmarks such as SciFact \citep{wadden2020scifact} learn to retrieve and rationalise sentence-level evidence for or against a claim before issuing a verdict. Compared to Hypothetica, OpenNovelty is more claim-refutation-oriented and more centred on evidence verification and taxonomy construction, whereas Hypothetica is more directly focused on user-facing originality analysis through four-dimensional scoring and sentence-level overlap mapping over a heterogeneous evidence base that combines academic works, patents, and code.

\citet{dasilva2025graphmind} introduce GraphMind, an interactive novelty-assessment tool that combines arXiv and Semantic Scholar APIs with LLM-based extraction to produce structured macro-level (paper graph) and micro-level (claim and section) views of related work, with explicit emphasis on traceability through an information-retrieval module that returns supporting and contrasting evidence. GraphMind is the closest existing tool to Hypothetica in terms of user-facing scope and evidence-traceability emphasis. The two systems differ in evidence base, with GraphMind drawing only from academic literature whereas Hypothetica retrieves additionally from patents and code through its adapter layer; in scoring layer, with GraphMind producing a qualitative report whereas Hypothetica produces a four-criterion Likert-based numerical score with categorical guardrails; and in interaction style, with GraphMind providing an exploratory web tool whereas Hypothetica runs an end-to-end pipeline that begins with a conversational interview to clarify the proposed idea before retrieval.

Beyond pure novelty checking, recent work has explored broader research-assistant systems built on LLM agents. \citet{eger2025survey} survey this fast-moving ecosystem across the full scientific lifecycle, including search, ideation, experimentation, content generation, and evaluation, providing a useful map for situating originality-assessment work among related AI-for-science tools. Within that ecosystem, SciMON \citep{wang2024scimon} explicitly optimises for novelty during scientific hypothesis generation by retrieving and contrasting against prior literature, reinforcing the view that novelty must be measured against retrieved evidence rather than assumed from the model alone. \citet{si2024llmideas} conduct a large-scale human study on whether LLMs can generate novel research ideas, finding that LLM-produced ideas are sometimes perceived as more novel than human-produced ideas but weaker in feasibility, which underscores the importance of grounding any novelty signal in retrieved evidence. ResearchAgent \citep{baek2024researchagent} iteratively refines research ideas using literature-grounded LLM agents, while The AI Scientist \citep{lu2024aiscientist} proposes an end-to-end pipeline for automated open-ended scientific discovery covering idea generation, experiment design, and write-up. These systems share Hypothetica's reliance on coordinated LLM agents but pursue idea generation, refinement, or experimentation rather than originality assessment of a user-supplied idea.

\subsection{Dense Retrieval and Cross-Encoder Reranking}

Established information-retrieval practice combines sparse lexical matching, dense embedding retrieval, and a learned reranker. Sparse methods such as BM25 \citep{robertson2009bm25} remain strong baselines for short queries with rare terms, and the BEIR benchmark \citep{thakur2021beir} shows that dense retrievers do not uniformly dominate sparse ones in zero-shot settings, which motivates the use of learned rerankers on top of dense first-stage retrieval. Dense retrievers project queries and documents into a shared vector space; the E5 family of contrastively pre-trained encoders \citep{wang2022e5} is one such family, used in this work. Cross-encoder rerankers, such as the BERT-based passage re-ranker introduced by \citet{nogueira2019monobert}, jointly encode query and document and produce relevance scores that are usually more accurate than pure bi-encoder similarity, at the cost of higher inference time.

Hypothetica adopts the standard two-stage design. It uses E5 embeddings stored in Chroma \citep{chroma} for first-stage retrieval, then applies cross-encoder reranking over a top-k candidate set before any document is forwarded to the scoring layer. For the OpenAlex adapter, an additional pre-ranking step weighs first-stage relevance, citation count, and recency before the embedding stage, because the unfiltered OpenAlex result set is much larger than the arXiv result set and benefits from a coarse pre-filter before any expensive computation.

\subsection{LLM-as-Judge and Rubric-Based Evaluation}

Hypothetica's per-paper scoring layer (Layer~1) follows the LLM-as-judge paradigm, in which a language model rates a target item against a set of criteria. G-Eval \citep{liu2023geval} proposed using GPT-4 with chain-of-thought prompting and form-filling to evaluate generated text against task-specific rubrics, reporting better correlation with human judgements than earlier reference-based metrics. MT-Bench and Chatbot Arena \citep{zheng2023mtbench} systematically studied LLM-as-judge agreement with human ratings on multi-turn benchmarks and identified well-known biases, including position bias and verbosity bias. Prometheus \citep{kim2024prometheus} showed that fine-grained, rubric-based evaluation does not necessarily require a frontier closed model: an open model trained on rubric-supervised data can match GPT-4 quality on many evaluation tasks. Closer to our setting, \citet{liang2024llmfeedback} conducted a large-scale empirical analysis of LLM-generated feedback on research papers across multiple journals and a major machine-learning conference, reporting that LLM feedback overlaps substantially with human reviewer comments while tending to be less specific and more uniform than human reviews. That study targets free-form reviewer-style feedback rather than per-criterion originality scoring, but it confirms that LLMs can produce useful, partially human-aligned judgements over academic text when prompted carefully.

These works inform two specific choices in Hypothetica. Each criterion is scored independently with its own Likert rubric in order to reduce halo effects, and the criterion weights used in Layer~2 are not uniform but calibrated empirically against a labelled arXiv benchmark, so that criteria with stronger discriminative power dominate the global score. Categorical guardrails additionally constrain the global score whenever any single criterion reaches its top Likert level, which mitigates the verbosity and overconfidence biases that pure averaging would otherwise allow.

\subsection{Multi-Agent LLM Architectures}

Agentic decomposition has been adopted to make complex LLM pipelines more structured and easier to audit. Surveys of LLM-based multi-agent systems describe common architectures and recurring issues such as coordination, verification, and reliability \citep{guo2024multiagent,zheng2025}. Frameworks such as AutoGen \citep{wu2023autogen} formalise multi-agent coordination among LLM agents through structured message passing, and broader surveys of retrieval-augmented generation \citep{gao2024ragsurvey} argue for similarly modular pipelines that separate retrieval, generation, and verification.

Hypothetica follows this trend by dividing the pipeline into specialised agents: a conversational interview agent that clarifies the proposed idea over a short multi-round dialogue, a reality-check agent that runs in parallel to flag potentially already-existing products, a query-variant agent that generates adapter-aware search queries, a relevant-paper-selector agent that filters reranked candidates, a Layer~1 per-paper analysis agent that scores the four originality criteria with per-criterion Likert rubrics, and a Layer~2 aggregation step that combines per-paper scores into a global originality score with categorical guardrails. For GitHub-based evidence, additional repository-relevance and synthesis agents replace the chunk-based retrieval flow. The pipeline requires structured outputs at every stage so that downstream scoring stays adequately deterministic and the user interface can provide clear explanations for each score.

\subsection{Beyond Academic Literature: Patents and Code as Prior Art}

Academic preprints are not the only form of prior art relevant to research originality. Industrial novelty is documented in patent records, and working systems are documented in public code repositories. Both have received attention from the information-retrieval community.

Patent prior-art retrieval is a long-standing task with its own conventions and benchmarks. \citet{krestel2021patents} survey deep-learning approaches to patent analysis, including classification, retrieval, and similarity assessment, and emphasise that patents differ from academic abstracts in length, terminology, and legal-style framing. These differences justify including patents as a first-class evidence source in originality assessment, because a research idea that is already disclosed in a patent is not novel even if no academic paper has yet been published on it.

Public code repositories provide a third channel of prior art that academic-only systems miss. From a novelty-assessment standpoint, the existence of a maintained repository that implements the proposed idea is a direct, often stronger signal than a paper that only describes it. Hypothetica embeds repository README files as natural-language documents using the same E5 encoder used for paper abstracts, rather than relying on code-specific embeddings, on the grounds that READMEs are written in prose and the originality decision is made by the LLM scoring agent rather than by the embedding similarity alone.

Hypothetica integrates these two evidence types alongside academic works. Patents are searched through Google Patents \citep{serpapi}, and repositories are searched on GitHub \citep{githubAPI}. Both adapters share the same pipeline interface as the OpenAlex and arXiv adapters, so the same scoring agents reason uniformly over heterogeneous sources without source-specific adjustments.

\subsection{Positioning of Hypothetica}

Hypothetica uses an evidence-first method aligned with retrieval-augmented generation. Instead of returning a single opaque score, it reports sentence-level flags and points to matching excerpts in the retrieved sources, so users can check the evidence behind each flagged passage. The system is therefore useful not only for producing a final score but also for showing where overlap occurs and why. While the initial Hypothetica prototype drew its corpus solely from arXiv using the official arXiv API \citep{arxivAPI}, the current system retrieves evidence through a pluggable adapter layer that integrates OpenAlex \citep{openalexDocs,priem2022openalex} for broad cross-disciplinary scholarly works, arXiv \citep{arxivAPI} for preprints in computer science and machine learning, Google Patents for industrial prior art, and GitHub repositories for working open-source implementations whose existence often invalidates novelty claims even when no paper has yet been published. The same interface can accommodate additional sources such as the Semantic Scholar API \citep{semanticScholarAPI} without changes to the scoring agents. In this respect, Hypothetica aligns with the recent movement toward grounded, multi-source, and inspectable originality analysis, while remaining distinct from generic similarity tools and from purely parametric LLM judgements.
